\begin{theorem}[Finite deck quotients and central cycles]\hypertarget{tag-0301}{}\label{push:src:B.1.2}\label{push:B12}
\StatementRecord{push:src:B.1.2}{proved}{0301}{\SeriesRef{push:src:B.1.1}, \SeriesRef{fam:src:E.2.1}, \SeriesRef{fam:src:A.5.1}}

The twisted finite deck generators extend to the models of
\ref{push:src:B.1.1}. On the twelve component labels their actions are
$i\mapsto i+4$ at $j=1$ and $i\mapsto i-3$ at $j=2$,
with labels in $\mathbb Z/12$. Every nonidentity group
element acts freely, including at every node of the central
cycle. The quotient is a smooth complex fourfold, proper
and flat over $(t,r)=(s_j^{m_j},r)$, with a reduced
normal-crossings $r$-central divisor. Its fibre divisor
over $t=p_j$ has multiplicity exactly $m_j$.

The reduction over $s_j=r=0$ is the central surface of \cref{fam:src:E.2.1}:
a cycle of four ruled surfaces at $j=1$, and a cycle of
three at $j=2$. In the root presentation the normal
character to $s_j=0$ is $e^{-2\pi i/m_j}$ and the
character on $r$ is trivial.
\end{theorem}


\begin{proof}
\begin{enumerate}[label=\arabic*.,leftmargin=*]
\item We retain source and target arguments. Write primes for
evaluation at $g_js$, put $\zeta_j=e^{-2\pi i/m_j}$,
and set
\[
 R_1=\begin{pmatrix}-1/\tau&0\\(1-\mu)/\tau&1\end{pmatrix},
 \qquad
 R_2=\begin{pmatrix}1/\tau&0\\-\mu/\tau&1\end{pmatrix}.
\]
The identities $R_j\Pi_c=\Pi'_c A_j$ and $A_jv_j=v_j$
follow by substitution of the displayed matrices. In the coordinates
$(s,\xi,W,r)$ the two affine deck maps become
\begin{align*}
 s'&=\zeta_1s,\qquad
 \xi'=-\xi/\tau+(6\mu'+2\tau'-4)/3,\\
 W'&=r^4 U_1(s,\xi)W,\\
 U_1&=\exp\!\left[-2\pi i
  \left(\frac{1-\mu}{\tau}\xi+
        \frac{\beta^{{\rm part}\prime}+2\mu'}3\right)\right],
 \end{align*}
and
\begin{align*}
 s'&=\zeta_2s,\qquad
 \xi'=\xi/\tau+(-6\mu'-3\tau'+3)/4,\\
 W'&=r^{-3} U_2(s,\xi)W,\\
 U_2&=\exp\!\left[-2\pi i
  \left(-\frac{\mu}{\tau}\xi-
        \frac{\beta^{{\rm part}\prime}+3\mu'}4\right)\right].
 \end{align*}
Each $U_j$ is a nowhere vanishing holomorphic function on
the additive elliptic cover. Period covariance gives the descent rule for $U_j$ and
therefore the required bundle isomorphism after quotient
by $\ker\gamma$. Multiplication by a unit preserves every toric chart.
Multiplication by $r^a$ sends the chart of component index
$i$ to the chart of index $i+a$. Therefore the displayed
maps and their inverses extend holomorphically to the
resolved infinite chain and then to its twelve-component
quotient. The identities $g_1^3=g_2^4=1$ on the quotient
hold on the dense punctured family, and hence everywhere.

\item A nonidentity power has component shift $4,8$ in the first
case and $-3,-6,-9$ in the second. None is zero modulo
twelve. The same shifts act on the edge labels of the
cycle, so no node can be fixed either. At $s\ne0$ the
rotation of $s$ already excludes fixed points. The finite
quotient map is consequently a holomorphic local
isomorphism. Its source is smooth, so its target is smooth.
Properness follows from the proper root model and the
finite base map. Flatness can be checked after the
surjective local-isomorphism cover: there $s^{m_j}$ and
$r$ are a regular sequence, and their fibres have pure
dimension two. The local central equation remains
$r=XY$ or $r=X$. Finally $t=s^{m_j}$ shows that the
Cartier divisor over $p_j$ is exactly $m_j$ times its
reduction, including where this divisor meets a central
node. This also proves the asserted normal characters.

\item At $s=0$ use the explicitly marked isogenies
\[
 L_1(\xi,\eta)=(\eta,\xi-\eta/3),\qquad
 L_2(\xi,\eta)=(\eta,\xi-\eta/2).
\]
Their kernel generators are $(2/3,\rho)$ and $(1/2,i)$,
and the finite deck maps lift to
\[
 (\xi,\eta)\mapsto(\xi+\kappa_1/3,(\rho-1)\eta),
 \qquad
 (\xi,\eta)\mapsto(\xi-\kappa_2/4,-i\eta).
\]
In the Tate coordinate $v=e^{-2\pi i\xi}$ these are the
shifts $4,-3$ just computed. The kernel generators act
by nonzero translations on the smooth elliptic factor
$F_j$, hence freely even at the two fixed points of each
Tate component. Every other element with nonzero deck
power has a nonzero component shift. The full product
groups are therefore free on the product semistable
models. Their component orbits have cardinalities three
and four, giving precisely four and three ruled
components. The original lattice isogenies intertwine
the maps on the dense tori and their unit-interval fans,
so the quotient identifications extend across all nodes.
This verifies the particular comparison with \cref{fam:src:E.2.1} that
we use.

\item For the punctured family obtained by descending the
$c$-family, the holomorphic correction $\eta_j(s)$ in
\cref{fam:src:A.5.1} causes no additional local obstruction here. At the
shift $c\mapsto c+12$, the reference translations
$-4\Pi_c(\widehat u+\widehat w)$ and
$3\Pi_c(\widehat u+\widehat w)$ are periods.
The remaining correction is equivariant and independent
of $c$. Translation by $-c\eta_j(s)/12$ on the additive
cover conjugates that descent to the period descent.
It commutes with the finite deck action by equivariance.
Thus the constructed finite models fill the actual local
pullback by $w=r^6$ after a specified holomorphic
identification over $r\ne0$. 
\end{enumerate}
\end{proof}
